\chapter{布尔函数与语言}
\label{chap:boolean-functions-and-languages}
\lecture
\label{lec:boolean-functions-and-languages}

\section{从一般函数到布尔函数}

一般的计算问题是函数

\[
  f:\bitsstar\longrightarrow\bitsstar.
\]

事实上，只研究取值于 $\bits$ 的布尔函数并不会损失一般性。给定 $f$，
定义它的逐位查询函数

\[
  b_f:\bitsstar\times\Nat\times\bits\longrightarrow\bits
\]

如下：

\begin{equation}
  \label{eq:bit-query-function}
  b_f(x,i,c)=
  \begin{cases}
    f(x)_i, & i<\card{f(x)}\text{ 且 }c=0,\\
    1,      & i<\card{f(x)}\text{ 且 }c=1,\\
    0,      & i\geq\card{f(x)}.
  \end{cases}
\end{equation}

这里从 $0$ 开始给输出位编号。自然数 $i$ 和三元组 $(x,i,c)$ 均使用固定的
二进制编码；具体选用哪一种合理编码，不影响下面的可计算性结论。

\begin{lemma}[逐位查询与原函数等价]
  \label{lem:bit-query-computability}
  函数 $f$ 可计算，当且仅当 $b_f$ 可计算。
\end{lemma}

\begin{proof}
  若有计算 $f$ 的程序，先算出 $y=f(x)$，再按
  \cref{eq:bit-query-function} 回答查询，即可计算 $b_f$。

  反过来，注意到

  \[
    b_f(x,i,1)=1\iff i<\card{f(x)},
    \qquad
    f(x)_i=b_f(x,i,0)\quad(i<\card{f(x)}).
  \]

  从 $i=0$ 开始依次查询 $b_f(x,i,1)$。只要结果为 $1$，就把
  $b_f(x,i,0)$ 追加到输出；第一次得到 $0$ 时停止。所得字符串恰为 $f(x)$。
\end{proof}

因此，讨论一般字符串函数的可计算性，可以归约为讨论布尔函数的可计算性。

\section{语言与布尔函数}

\begin{definition}[语言与示性函数]
  \label{def:language-characteristic-function}
  一个\emph{语言}（language）是二进制串集合 $A\subseteq\bitsstar$。
  对布尔函数 $f:\bitsstar\to\bits$，定义

  \[
    A_f=\{x\in\bitsstar:f(x)=1\}.
  \]

  反过来，语言 $A$ 的\emph{示性函数}（characteristic function）定义为

  \[
    \chi_A(x)=
    \begin{cases}
      1, & x\in A,\\
      0, & x\notin A.
    \end{cases}
  \]
\end{definition}

上述两个变换互为逆映射。因此布尔函数与语言一一对应，并且

\[
  f(x)=1\iff x\in A_f.
\]
